% GENERATED FILE. Edit canonical lessons, then run npm run generate:books.
% canonical-source-hash: fnv1a64:c14c47e535929a0f
% canonical-lessons: SA-C507-bhastra, SA-C507-sandamsah, SA-C507-vilapayati, SA-C507-dhamati, SA-C507-bhusanam

\chapter{Metal and Gold}
\label{ch:sa-metal-and-gold}
\hlchaptermodality{507}

\begin{chapteropening}
\textbf{By the end of this chapter:} \emph{I can describe how a smith melts, shapes and polishes metal.}

Complete the last lesson of chapter 507: I can describe how a smith melts, shapes and polishes metal.
\end{chapteropening}

\section[\texorpdfstring{\sk{भस्त्रा} (bhastrā)}{भस्त्रा (bhastrā)}]{\sk{भस्त्रा} (bhastrā) — a bellows}

\label{lesson:SA-C507-bhastra}



\begin{quote}
Before the new one: say the Sanskrit for dyes, then the Sanskrit for a tailor.
\end{quote}

\subsection*{You'll want to know: \sk{भस्त्रा}}
\textbf{\sk{भस्त्रा}} — \emph{bhastrā} — \textquotedblleft{}a bellows\textquotedblright{}.

\begin{quote}
\sk{लोहकारः} \sk{भस्त्रां} \sk{चालयति।}
\end{quote}

\begin{quote}
\emph{lohakāraḥ bhastrāṃ cālayati.}
\end{quote}

\begin{quote}
The blacksmith works the bellows.
\end{quote}

A feminine noun.

\begin{cousinweb}
\textbf{\sk{भस्त्रा}} (\emph{bhastrā}) is the word MW records for a leather bottle or bag, and so a bellows; MW gives no derivation, so its origin is uncertain (MW).
\end{cousinweb}

\begin{grammarlens}[title={What you've built}]
One more word for describing what you lived and felt.
\end{grammarlens}

\subsection*{Your turn}
\begin{itemize}
\raggedright
  \item \emph{Say it:} \emph{bhastrā}
  \item \emph{Say it:} \emph{bhastrā}, once more
  \item \emph{Say it:} \emph{lohakāraḥ bhastrāṃ cālayati.} — The blacksmith works the bellows.
  \item \emph{Recall:} say the Sanskrit for dyes, then the Sanskrit for a tailor, then say \emph{bhastrā} again
\end{itemize}

\begin{tcolorbox}[breakable,colback=teal!4,colframe=teal!35!black,title={Before you move on}]
What does \sk{भस्त्रा} mean? (\textquotedblleft{}A bellows\textquotedblright{}.) Say it once more.
\end{tcolorbox}
\section[\texorpdfstring{\sk{सन्दंशः} (sandaṃśaḥ)}{सन्दंशः (sandaṃśaḥ)}]{\sk{सन्दंशः} (sandaṃśaḥ) — tongs}

\label{lesson:SA-C507-sandamsah}



\begin{quote}
Before the new one: say the Sanskrit for a tailor, then the Sanskrit for a bellows.
\end{quote}

\subsection*{You'll want to know: \sk{सन्दंशः}}
\textbf{\sk{सन्दंशः}} — \emph{sandaṃśaḥ} — \textquotedblleft{}tongs\textquotedblright{}.

\begin{quote}
\sk{लोहकारः} \sk{सन्दंशेन} \sk{अंगारं} \sk{गृह्णाति।}
\end{quote}

\begin{quote}
\emph{lohakāraḥ sandaṃśena aṅgāraṃ gṛhṇāti.}
\end{quote}

\begin{quote}
The blacksmith picks up a live coal with the tongs.
\end{quote}

A masculine noun.

\begin{cousinweb}
\textbf{\sk{सन्दंशः}} (\emph{sandaṃśaḥ}) is built on \emph{sam}, \textquotedblleft{}together\textquotedblright{}, and the root √daṃś, \textquotedblleft{}to bite\textquotedblright{}: a biting together; MW gives \emph{saṃdaṃśa} as a pair of tongs or pincers (MW).
\end{cousinweb}

\begin{grammarlens}[title={What you've built}]
One more word for describing what you lived and felt.
\end{grammarlens}

\subsection*{Your turn}
\begin{itemize}
\raggedright
  \item \emph{Say it:} \emph{sandaṃśaḥ}
  \item \emph{Say it:} \emph{sandaṃśaḥ}, once more
  \item \emph{Say it:} \emph{lohakāraḥ sandaṃśena aṅgāraṃ gṛhṇāti.} — The blacksmith picks up a live coal with the tongs.
  \item \emph{Recall:} say the Sanskrit for a tailor, then the Sanskrit for a bellows, then say \emph{sandaṃśaḥ} again
\end{itemize}

\begin{tcolorbox}[breakable,colback=teal!4,colframe=teal!35!black,title={Before you move on}]
What does \sk{सन्दंशः} mean? (\textquotedblleft{}Tongs\textquotedblright{}.) Say it once more.
\end{tcolorbox}
\section[\texorpdfstring{\sk{विलापयति} (vilāpayati)}{विलापयति (vilāpayati)}]{\sk{विलापयति} (vilāpayati) — melts down}

\label{lesson:SA-C507-vilapayati}



\begin{quote}
Before the new one: say the Sanskrit for a bellows, then the Sanskrit for tongs.
\end{quote}

\subsection*{You'll want to know: \sk{विलापयति}}
\textbf{\sk{विलापयति}} — \emph{vilāpayati} — \textquotedblleft{}melts down\textquotedblright{}.

\begin{quote}
\sk{स्वर्णकारः} \sk{अग्नौ} \sk{सुवर्णं} \sk{विलापयति।}
\end{quote}

\begin{quote}
\emph{svarṇakāraḥ agnau suvarṇaṃ vilāpayati.}
\end{quote}

\begin{quote}
The goldsmith melts the gold in the fire.
\end{quote}

Past (he …): \sk{सः} \sk{सुवर्णं} \sk{विलापितवान्}, \emph{saḥ suvarṇaṃ vilāpitavān}.

\begin{cousinweb}
\textbf{\sk{विलापयति}} (\emph{vilāpayati}) is built on \emph{vi}, \textquotedblleft{}apart\textquotedblright{}, and the causative of the root √lī, \textquotedblleft{}to cling, to dissolve\textquotedblright{}; MW gives \emph{vilī} the sense \textquotedblleft{}to melt\textquotedblright{}, and its causative \emph{vilāpayati} makes a thing melt (MW).
\end{cousinweb}

\begin{grammarlens}[title={What you've built}]
One more word for describing what you lived and felt.
\end{grammarlens}

\subsection*{Your turn}
\begin{itemize}
\raggedright
  \item \emph{Say it:} \emph{vilāpayati}
  \item \emph{Say it:} \emph{vilāpayati}, once more
  \item \emph{Say it:} \emph{svarṇakāraḥ agnau suvarṇaṃ vilāpayati.} — The goldsmith melts the gold in the fire.
  \item \emph{Recall:} say the Sanskrit for a bellows, then the Sanskrit for tongs, then say \emph{vilāpayati} again
\end{itemize}

\begin{tcolorbox}[breakable,colback=teal!4,colframe=teal!35!black,title={Before you move on}]
What does \sk{विलापयति} mean? (\textquotedblleft{}Melts down\textquotedblright{}.) Say it once more.
\end{tcolorbox}
\section[\texorpdfstring{\sk{धमति} (dhamati)}{धमति (dhamati)}]{\sk{धमति} (dhamati) — blows on a fire}

\label{lesson:SA-C507-dhamati}



\begin{quote}
Before the new one: say the Sanskrit for tongs, then the Sanskrit for melts down.
\end{quote}

\subsection*{You'll want to know: \sk{धमति}}
\textbf{\sk{धमति}} — \emph{dhamati} — \textquotedblleft{}blows on a fire\textquotedblright{}.

\begin{quote}
\sk{लोहकारः} \sk{भस्त्रया} \sk{अग्निं} \sk{धमति।}
\end{quote}

\begin{quote}
\emph{lohakāraḥ bhastrayā agniṃ dhamati.}
\end{quote}

\begin{quote}
The blacksmith blows on the fire with the bellows.
\end{quote}

Past (he …): \sk{सः} \sk{शंखं} \sk{ध्मातवान्}, \emph{saḥ śaṅkhaṃ dhmātavān}.

\begin{cousinweb}
\textbf{\sk{धमति}} (\emph{dhamati}) comes from the root √dham (or \emph{dhmā}), \textquotedblleft{}to blow\textquotedblright{}, as one blows a fire, a conch or a bellows; its past participle is \emph{dhmāta} (MW).
\end{cousinweb}

\begin{grammarlens}[title={What you've built}]
One more word for describing what you lived and felt.
\end{grammarlens}

\subsection*{Your turn}
\begin{itemize}
\raggedright
  \item \emph{Say it:} \emph{dhamati}
  \item \emph{Say it:} \emph{dhamati}, once more
  \item \emph{Say it:} \emph{lohakāraḥ bhastrayā agniṃ dhamati.} — The blacksmith blows on the fire with the bellows.
  \item \emph{Recall:} say the Sanskrit for tongs, then the Sanskrit for melts down, then say \emph{dhamati} again
\end{itemize}

\begin{tcolorbox}[breakable,colback=teal!4,colframe=teal!35!black,title={Before you move on}]
What does \sk{धमति} mean? (\textquotedblleft{}Blows on a fire\textquotedblright{}.) Say it once more.
\end{tcolorbox}
\section[\texorpdfstring{\sk{भूषणम्} (bhūṣaṇam)}{भूषणम् (bhūṣaṇam)}]{\sk{भूषणम्} (bhūṣaṇam) — an ornament}

\label{lesson:SA-C507-bhusanam}



\begin{quote}
Before the new one: say the Sanskrit for melts down, then the Sanskrit for blows on a fire.
\end{quote}

\subsection*{You'll want to know: \sk{भूषणम्}}
\textbf{\sk{भूषणम्}} — \emph{bhūṣaṇam} — \textquotedblleft{}an ornament\textquotedblright{}.

\begin{quote}
\sk{स्वर्णकारः} \sk{सुवर्णात्} \sk{भूषणानि} \sk{रचयति।}
\end{quote}

\begin{quote}
\emph{svarṇakāraḥ suvarṇāt bhūṣaṇāni racayati.}
\end{quote}

\begin{quote}
The goldsmith makes ornaments from gold.
\end{quote}

A neuter noun.

\begin{cousinweb}
\textbf{\sk{भूषणम्}} (\emph{bhūṣaṇam}) comes from the root √bhūṣ, \textquotedblleft{}to adorn\textquotedblright{}; MW gives \emph{bhūṣaṇa} as decorating, an ornament (MW).
\end{cousinweb}

\begin{grammarlens}[title={What you've built}]
One more word for describing what you lived and felt.
\end{grammarlens}

\subsection*{Your turn}
\begin{itemize}
\raggedright
  \item \emph{Say it:} \emph{bhūṣaṇam}
  \item \emph{Say it:} \emph{bhūṣaṇam}, once more
  \item \emph{Say it:} \emph{svarṇakāraḥ suvarṇāt bhūṣaṇāni racayati.} — The goldsmith makes ornaments from gold.
  \item \emph{Recall:} say the Sanskrit for melts down, then the Sanskrit for blows on a fire, then say \emph{bhūṣaṇam} again
\end{itemize}

\begin{tcolorbox}[breakable,colback=teal!4,colframe=teal!35!black,title={Before you move on}]
What does \sk{भूषणम्} mean? (\textquotedblleft{}An ornament\textquotedblright{}.) Say it once more.
\end{tcolorbox}
